\documentclass[11pt]{article}
\usepackage[margin=1in]{geometry}
\usepackage{amsmath,amssymb,amsthm,booktabs,tabularx,array,enumitem,hyperref}
\newtheorem{definition}{Definition}
\newtheorem{proposition}{Proposition}
\newcommand{\Met}{\texttt{met}}
\newcommand{\NotMet}{\texttt{not\_met}}
\newcommand{\Indet}{\texttt{indeterminate}}
\newcommand{\AC}{\mathsf{AC}}
\newcommand{\UN}{\mathsf{UN}}
\newcommand{\IV}{\mathsf{IV}}
\newcommand{\ED}{\mathsf{ED}}
\newcommand{\DS}{\mathsf{DS}}
\newcommand{\MQ}{\mathsf{MQ}}
\newcommand{\PF}{\mathsf{PF}}

\title{Non-Empirical Ontology Success Criteria\\
\large P14-T02 task-local draft/control formalization}
\author{The AEther-Flow Research Project}
\date{22 July 2026}

\begin{document}
\maketitle

\begin{abstract}
This packet defines seven independent, non-compensatory evidence criteria for
assessing an ontology that is empirically equivalent to a benchmark:
assumption compression, unification, inevitability, explanatory depth,
dimension-and-signature explanation, matter-or-quantum connection, and
problem-solving fertility.  Each coordinate requires its own evidence
certificate and is reported with a tri-valued status rather than a score.
Applied to the current AEther-Flow record, six obligations are not met and
problem-solving fertility remains indeterminate.  These are present evidence
statuses, not impossibility results.  The packet does not select or modify an
ontology, change scientific ledgers, promote a claim, or establish a
first-principles derivation of general relativity.
\end{abstract}

\section{Scope, comparator, and status algebra}

Let \(x\) be a precisely scoped ontology proposal, let \(B\) be a declared
benchmark such as ordinary general relativity on the exact-GR operational
domain, and let \(E(x)\) be the source-backed evidence admitted for \(x\).
Empirical equivalence to \(B\) fixes the operational comparison; it is not
itself evidence that \(x\) has additional non-empirical value.

\begin{definition}[Tri-valued evidence status]
For each criterion \(c\), define
\[
 s_c(x\mid B)\in
 S=\{\Met,\NotMet,\Indet\}.
\]
The value \(\Met\) requires a complete criterion-specific evidence
certificate.  The value \(\NotMet\) means that at least one declared
obligation is known to be unsatisfied in the present record; it is not a theorem of impossibility.
The value \(\Indet\) means that the available
source record does not decide whether the complete obligation is satisfied.
\end{definition}

\begin{definition}[Evidence certificate]
A certificate for criterion \(c\) is a tuple
\[
 \operatorname{Cert}_{c}(x\mid B)
   =(H_c,P_c,W_c,A_c,R_c),
\]
where \(H_c\) states the scope and hypotheses, \(P_c\) is an acyclic
source-to-result premise map, \(W_c\) is a criterion-specific witness or
theorem, \(A_c\) compares relevant alternatives and the benchmark, and
\(R_c\) records robustness and failure conditions.  A validator, registry
row, task status, artifact count, role label, or narrative assertion cannot
replace \(P_c\) or \(W_c\).
\end{definition}

The non-empirical evidence vector is
\[
 \Gamma(x\mid B)=
 (s_{\AC},s_{\UN},s_{\IV},s_{\ED},s_{\DS},s_{\MQ},s_{\PF})\in S^7.
\]
Its coordinates are not probabilities, utilities, or measurements of truth.

\section{Seven criterion families}

\begin{definition}[Assumption compression, \(\AC\)]
Assumption compression is met only when a common consequence set is recovered
from a demonstrably smaller or less independent source premise basis than the
benchmark-relative formulation, with an explicit translation and an audit
showing that removed premises do not return inside bridge maps, boundary
conditions, selectors, or reconstruction rules.
\end{definition}
An admissible example is an acyclic derivation that replaces several
independent postulates by one source law and proves the old postulates as
corollaries.  Renaming GR variables while adding a substrate vocabulary is a
counterexample: it can preserve every benchmark premise while increasing the
unproved source debt.

\begin{definition}[Unification, \(\UN\)]
Unification is met only when the same independently specified source
structure or law derives at least two previously separate result families,
including their compatibility relation, without merely placing them in one
document or notation.
\end{definition}
A single source dynamics that yields both geometrical response and matter
coupling under shared hypotheses would be relevant evidence.  A common
glossary, common role, common data format, or shared exact-GR action is not a
unification certificate.

\begin{definition}[Inevitability, \(\IV\)]
Inevitability is met only relative to a declared admissible alternative class.
The certificate must prove uniqueness up to a stated equivalence, or prove
robust dynamical selection with nonzero basin or measure under source-only
conditions.  Every selector, quotient, gauge, and boundary choice must be
identified.
\end{definition}
A uniqueness theorem over a source-defined candidate class is positive
evidence.  Choosing one representative, atlas, dimension, signature, or
normalization because it reproduces the target is a counterexample.

\begin{definition}[Explanatory depth, \(\ED\)]
Explanatory depth is met only when the source model supplies a non-circular
mechanism or dependency structure that answers contrastive and
counterfactual questions not answered by the benchmark description alone.
The certificate must name interventions or controlled variations, their
source-side consequences, and the limits of the explanation.
\end{definition}
Deriving why a stable target structure changes under a specified source
variation is relevant evidence.  Narrative coherence, intuitive imagery,
familiar mathematics, or a dictionary that leaves every benchmark dependency
unchanged is insufficient.

\begin{definition}[Dimension-and-signature explanation, \(\DS\)]
This criterion has two mandatory subcertificates.  The dimension certificate
must derive the effective dimension and exclude or control relevant
alternatives.  The signature certificate must derive Lorentzian signature,
time orientation, and the appropriate invariance class from source-only
structure.  Both must be stable under the declared variations and neither may
assume a four-dimensional smooth Lorentzian target.
\end{definition}
A source theorem selecting a four-dimensional Lorentzian effective sector
with stability bounds could satisfy the criterion.  Positing a
four-dimensional differentiable substrate, or importing the target metric and
then reading off its signature, does not.

\begin{definition}[Matter-or-quantum connection, \(\MQ\)]
Matter-or-quantum connection is met only when source primitives construct a
typed matter or quantum sector, its dynamics and observables, and a controlled
bridge to the benchmark limit.  If coupling is claimed, universality,
conservation, operational semantics, and robustness must be discharged in the
same declared scope.
\end{definition}
Deriving representations, coupling, and a tested low-energy limit from a
source law is relevant evidence.  Inserting the ordinary matter action,
stress-energy tensor, detector semantics, Hilbert space, or quantum rule by
hand is not.

\begin{definition}[Problem-solving fertility, \(\PF\)]
Problem-solving fertility is met only when the ontology enables an
independently checkable result on a predeclared problem: a new theorem,
constraint, countermodel, calculation, or empirical discriminator whose
dependence on the ontology is explicit and whose value is not reducible to
workflow output.
\end{definition}
A source principle that resolves a previously open construction problem and
survives independent checking is relevant evidence.  Counts of agents,
artifacts, candidates, validators, commits, or PASS receipts are
counterexamples.

\section{Non-compensation and comparison rule}

\begin{proposition}[NEO-NONCOMPENSATION-PROPOSITION-V1]
Under the certificate semantics above, no weighted sum of criterion labels is
a sufficient non-empirical success certificate.  In particular, a met
coordinate cannot discharge a missing witness, circular premise map, or
failed source boundary in another coordinate.
\end{proposition}

\begin{proof}
Each status \(\Met\) is defined by the existence of a typed,
criterion-specific certificate.  A weighted sum retains only assigned numeric
labels and weights; it does not retain the premise graph, witness, alternative
class, robustness conditions, or source authority.  Two records can therefore
receive the same sum while one contains a valid certificate and the other
contains a circular or target-importing surrogate.  Conversely, increasing a
different coordinate can raise the sum without supplying the missing
certificate.  The sum is therefore not sufficient for any coordinate and
cannot authorize success.  Reporting must remain coordinate-wise, with
uncertainty and vetoing defects explicit.
\end{proof}

Pairwise comparison is permitted only by displaying both evidence vectors,
their certificate scopes, and any strict coordinate-wise advantage supported
by source evidence.  Mixed vectors remain incomparable.  Aesthetic
preference, engineering convenience, or chosen weights may inform protected
human deliberation but are not scientific evidence and do not select an
ontology.

\section{Current source-grounded assessment}

The current exact-GR line is already classified as an interpretive
redescription, not a completed substrate derivation.  The comprehensive
assumption ledger records four-dimensionality and differentiability as
premises, leaves Lorentzian signature, physical scale, source matter
semantics, universal source coupling, physical covariance, and substrate
dynamics without derivation-critical source laws, and records the target GR
action as adopted benchmark closure.  The fixed continuum-first and
emergence-first options are proposal-only alternatives; the comparative audit
found neither scientifically dominant.

\begin{center}
\small
\begin{tabularx}{\textwidth}{>{\raggedright\arraybackslash}p{0.25\linewidth}
  >{\raggedright\arraybackslash}p{0.14\linewidth}
  >{\raggedright\arraybackslash}X}
\toprule
Criterion & Status & Present evidence boundary \\
\midrule
Assumption compression & \NotMet & No acyclic certificate shows fewer independent premises than the exact-GR baseline; key source and bridge obligations remain open. \\
Unification & \NotMet & A common substrate vocabulary is present, but no source law derives geometry and matter or quantum structure together. \\
Inevitability & \NotMet & Multiple proposal-only regimes and selector choices remain; no source-only uniqueness or robust-selection theorem chooses among them. \\
Explanatory depth & \NotMet & The interpretive dictionary is disciplined, but no source mechanism currently supplies the required counterfactual dependency certificate. \\
Dimension and signature & \NotMet & Four-dimensionality is posited in the canonical foundations, while Lorentzian signature belongs to adopted exact-GR closure rather than a source-only derivation. \\
Matter or quantum connection & \NotMet & Ordinary matter coupling is adopted in the benchmark; source matter semantics, coupling, and a quantum bridge remain open. \\
Problem-solving fertility & \Indet & The program has task-local mathematical results and obstructions, but their ontology-dependent, independently reusable problem-solving value has not been established. \\
\bottomrule
\end{tabularx}
\end{center}

Thus
\[
\Gamma(\text{current AEther-Flow}\mid\text{GR})
=(\NotMet,\NotMet,\NotMet,\NotMet,\NotMet,\NotMet,\Indet).
\]
This vector does not retract the established exact-GR interpretive
redescription, and it does not prove that future source extensions cannot
improve any coordinate.

\section{Material evidence that would improve the record}

\begin{enumerate}[leftmargin=*]
  \item \(\AC\): an acyclic premise map and comparison proof showing a genuine
  reduction of independent assumptions without hidden reintroduction.
  \item \(\UN\): one source law deriving at least two currently separate
  sectors and their compatibility conditions.
  \item \(\IV\): a source-defined uniqueness theorem or robust dynamical
  selection result over explicit alternatives.
  \item \(\ED\): a mechanism with controlled source interventions and novel
  contrastive consequences.
  \item \(\DS\): separate source-only dimension and Lorentzian-signature
  theorems with stability and non-target-import audits.
  \item \(\MQ\): source-derived matter or quantum structure with dynamics,
  observables, coupling where claimed, and a controlled benchmark limit.
  \item \(\PF\): an ontology-dependent result on a predeclared problem that is
  independently checked and not reducible to process output.
\end{enumerate}

\section{Authority boundary}

This is a task-local science-audit result with status
\texttt{draft/control}.  It does not edit or reinterpret canonical ontology,
select or reject the continuum-first or emergence-first proposal, adopt a
source law, change the Distance-to-GR ledger, establish physical gauge status,
derive matter or quantum physics, promote a benchmark, authorize publication,
issue a Gate Chair verdict, execute P14-T03, or claim a completed derivation.

\section*{Source materials}

The AEther-Flow Research Project. (2026a). \emph{Ontology: Foundations}
[Canonical ontology source].
\texttt{ontology/tex/aether\_flow\_foundations.tex}.

The AEther-Flow Research Project. (2026b). \emph{Comprehensive
source-assumption ledger v1} [Task-local draft/control artifact].
\texttt{research\_control/tasks/RT-20260720-030/artifacts/
comprehensive\_source\_assumption\_ledger\_v1.tex}.

The AEther-Flow Research Project. (2026c). \emph{P4-T04 comparative audit of
the continuum-first and emergence-first ontology options} [Science-audit
artifact].
\texttt{research\_control/tasks/RT-20260721-002/artifacts/
ontology\_regime\_comparison\_report\_v1.md}.

The AEther-Flow Research Project. (2026d). \emph{Four project-success
categories and their evidence obligations} [Task-local draft/control
formalization].
\texttt{research\_control/tasks/RT-20260722-016/artifacts/
four\_project\_success\_categories\_v1.tex}.

\end{document}
